%%%
\section{Context and Motivations\label{sec:intro}}

The success of Cloud Computing has driven the advent of Utility Computing. However, Cloud
Computing is a victim of its own success: In order to answer the escalating demand for
computing resources, Cloud Computing providers must build data centers~(DCs) of
ever-increasing size. As a consequence, besides facing the well-known issues of large-scale platforms
management, large-scale DCs have now to deal with energy considerations that limit the number
of physical resources that one location can host.

Instead of investigating alternative solutions that could tackle the aforementioned
concerns, the current trend consists in deploying larger and larger DCs in few strategic
locations presenting energy advantages. For example, Western North Carolina, USA, is an
attractive area due to its abundant capacity of coal and nuclear power brought about the
departure of the textile and furniture industry~\cite{greenpeace:2013}. More recently,
several proposals suggested building next generation DCs close to the polar circle in
order to leverage free cooling techniques, considering that cooling is accounting for a
big part of the electricity consumption~\cite{greenberg:sigcomm09}.

%%%
\subsection{Inherent Limitations of Large-scale DCs}

Although building large scale DCs  enables to cope with the actual demand,
% while continuing to operate UC resources through centralized software system 
it is far from delivering sustainable and efficient UC infrastructures. In addition to
requiring the construction and the deployment of a complete network infrastructure to
reach each DC, it exacerbates the inherent limitations of the Cloud Computing model:

\begin{itemize}
\item The externalization of private applications/data often faces legal issues that
  restrain companies from outsourcing them on external infrastructures, especially when
  located in other countries.
\item The overhead implied by the unavoidable use of the Internet to reach distant
  platforms is wasteful and costly in several situations: Deploying a broadcasting service
  of local events or an online service to order pizzas at the edge of the polar circle,
  for instance, leads to important overheads since most of the users are \emph{a priori}
  located in the neighborhood of the event/the pizzeria.
\item The connectivity to the application/data cannot be ensured by centralized dedicated
  centers, especially if they are located in a similar geographical zone. The only way to
  ensure disaster recovery is to leverage distinct sites.\footnote{``Amazon outages –
    lessons learned'',
    \href{http://gigaom.com/cloud/amazon-outages-lessons-learned/}{http://gigaom.com/cloud/amazon-outages-lessons-learned/}
    (valid on Nov 2013, the 30\textsuperscript{th}).}
\end{itemize}

The two first points could be partially tackled by hybrid or federated Cloud
solutions~\cite{armbrust:2010}, that aim at extending the resources available on one Cloud
with those of another one; however, the third one requires a disruptive change in the way
UC resources are managed.
%Deploying a local events broadcasting service or an
%online service to order pizza at the edge of the polar circle for instance, leads to an important overhead
%in terms of energy footprint, network exchanges as well as latency since it can be assumed
%that a vast majority of the users are located in the neighborhood of the event/the
%pizzeria.

Another issue is that, according to some projections of a recent IEEE
report~\cite{ieeenetreport:2012}, the network traffic continues to double roughly every
year. Consequently, bringing the IT services closer to the end-users is becoming crucial
to limit the energy impact of these exchanges and to save the bandwidth of some
links. Similarly, this notion of locality is critical for the adoption of the UC
model by applications that need to deal with a large amount of data as getting them in and
out actual UC infrastructures may significantly impact the global
performance~\cite{Fos11}.

The concept of micro/nano DCs at the edge of the backbone~\cite{greenberg:sigcomm09} may
be seen as a complementary solution to hybrid platforms in order to reduce the overhead of
network exchanges. However, operating multiple small DCs breaks somehow the idea of
mutualization in terms of physical resources and administration simplicity, making this
approach questionable.
% Moreover, the number of such micro/nano DCs will remain limited and the question of
% where and how federating a large number of such facilities are still not solved.

%%%
\subsection{Ubiquitous and Oversized Network Backbones}

One way to partially solve the mutualization concern enlightened by the defenders of
large-scale DCs is to directly deploy the concept of micro/nano DCs upon the Internet
backbone. People are (and will be) more and more surrounded by computing resources,
especially those in charge of interconnecting all IT equipments. Even though these small
and medium-sized facilities include resources that are barely
used~\cite{Andrew:2003,Benson:2010}, they can hardly be removed (\textit{e.g.} routers).
Considering this important aspect, we claim that a new generation of UC platforms can be
delivered by leveraging existing network centers, starting from the core nodes of the
backbone to the different network access points in charge of interconnecting public and
private institutions. By such a mean, network and UC providers would be able to mutualize
resources that are mandatory to operate network/data centers while delivering widely
distributed UC platforms able to better match the geographical dispersal of users.
%
% TODO: AL -> ACO, please introduce this point latter in the chapter
% As a consequence, several initiatives started investigating how they could be
%better leveraged to support the requirements and constraints of current IT
%usages.  The concept of \emph{data furnaces} \cite{liu:hotcloud11} is one of
%the promising idea that seeks to mitigate the cost of operating
%network/computing resources by using them as a source of heat inside public
%buildings such as hospitals or universities. 
%
Figure~\ref{fig:renater} allows to better capture the advantages of such a proposal.
%\ftodo[FQ$\rightarrow$ALL]{Nous n'avons pas encore introduit la contribution à ce stade $\rightarrow$ enlever cette phrase~?}
% We did, cf. above : WE CLAIM
It shows a snapshot of the network weather map of
RENATER\footnote{\href{http://www.renater.fr}{http://www.renater.fr}}, the backbone
dedicated to universities and research institutes in France. It reveals several
important points:
\begin{figure}[b]
\vspace*{-.3cm}
\includegraphics[width=10cm]{./FIGS/renater.png}
\centering\caption{The RENATER Weather Map on May 2013, the 27th, around 4PM.
Each red square corresponds to a particular point of presence (PoP) of the network. The map is available in real-time
at: \href{http://www.renater.fr/raccourci}{http://www.renater.fr/raccourci}}
\label{fig:renater}
\vspace*{-.3cm}
\end{figure}



\begin{itemize} 
\item As mentioned before, most of the resources are underutilized (only two links are used
  between 45\% and 55\%, a few between 25\% and 40\%, and the majority below the threshold
  of 25\%).
\item The backbone was deployed and is renewed to match the demand: The density of points
  of presence~(PoP, \ie a small or medium-sized network center) as well as the bandwidth
  of each link are more important on the edge of large cities such as Paris, Lyon or
  Marseille.
\item The backbone was designed to avoid disconnections, since 95\% of the PoPs can be
  reached by at least two distinct routes.
\end{itemize}


%\ftodo[AL -> ALL]{It might make sense to talk of congestion effects that we can see in Marseille and Paris, the two path to the rest of Internet.}

\subsection{Locality-based Utility Computing}

%This chapter aims at introducing a new generation of UC platforms that can be
%seen somehow as an extension of the concept of micro DCs. The main change is to
%consider locality as a key point of UC services and to leverage facilities
%composing the internet backbone.  %Instead of building and deploying dedicated
%facilities, we claim that next UC
%infrastructures should be tightly coupled with any facilities available through
%the Internet, starting from the core routers of the backbone, the different
%network access points and any small and medium-size computing infrastructures
%that may be provisioned by public and private institutions. 

% Although it involves radical changes in the way
%physical and virtual resources are managed, locating and operating computing
%power and data on
%facilities close to the end-users will deliver highly efficient
%and sustainable UC services, resolving inherent limitations of the cloud computing model leveraging large-scale DCs. 

This chapter aims at introducing locality-based UC infrastructures, a new generation of UC
platforms that solves inherent limitations of the Cloud Computing paradigm relying on
large-scale DCs. Although it involves radical changes in the way physical and virtual
resources are managed, leveraging network centers is a promising way to deliver highly
efficient and sustainable UC services.

From the physical point of view, network backbones 
%such as National Research and Educational Networks (NRENs) 
provide appropriate infrastructures, \ie, reliable and efficient enough to operate UC
resources spread across the different PoPs. Ideally, UC resources would be able to
directly take advantage of computation cycles available on network active devices,
\textit{i.e.} those in charge of routing packets. However, leveraging network resources to
make external computations may lead to important security concerns. Hence, we propose to
extend each PoP with a number of servers dedicated to VMs hosting. Because it is natural
to assume that the network traffic and UC demands are proportional, larger network centers
will be completed by more UC resources than the smaller ones. Moreover, by deploying UC
services on relevant PoPs, a LUC infrastructure will be able to natively confine network
exchanges to a minimal scope, minimizing altogether the energy footprint of the network, the
impact on latency and the congestion phenomena that may occur on critical paths (for
instance Paris and Marseille on RENATER).

From the software point of view, the main challenge is to design a comprehensive distributed
system in charge of turning a complex and diverse network of resources into a collection
of abstracted computing facilities that is both reliable and easy to operate.

\begin{svgraybox}
  The design of the \emph{LUC  Operating System},  an advanced system being able to unify many UC
  resources distributed on distinct sites,  would enable Internet service providers~(ISPs)
  and other institutions in charge of operating a network backbone to build an
  extreme-scale LUC infrastructure with a limited additional cost. Instead of redeploying
  a complete installation, they will be able to leverage IT resources and specific devices
  such as computer room air conditioning units, inverters or redundant power supplies
  already present in each center of their backbone.
\end{svgraybox}


\medskip
%This chapter  describes  how such a new
%generation of highly efficient and sustainable UC can emerge through an integrated
%system, \ie the \emph{LUC Operating System}, leveraging advanced and P2P system mechanisms.

In addition to considering \emph{locality} as a primary concern, the novelty of the LUC OS
proposal is to consider the VM as the basic object it manipulates.  Unlike existing
research on distributed operating systems designed around the process concept, a LUC OS
will manipulate VMs throughout a federation of widely distributed physical
machines. Virtualization technologies abstract out hardware heterogeneity, and allow
transparent deployment, preemption, and migration of virtual environments~(VEs), \ie a set
of interconnected VMs.  By dramatically increasing the flexibility of resource management,
virtualization allows to leverage state-of-the-art results from other distributed systems
areas such as autonomous and decentralized systems.  Our goal is to build a system that
allows end-users to launch VEs over a distributed infrastructure as simply as they launch
processes on a local machine, \ie without the burden of dealing with resources
availability or location.

\subsection{Chapter Outline} 
Section~\ref{sec:challenges} describes the key objectives of a LUC OS and the associated
challenges.  Section~\ref{sec:background} explains why our vision differs from current and
previous UC solutions. In Section~\ref{sec:archi}, we present how such a unified system
may be designed by delivering the premises of the \discovery system, an agent-based system
enabling distributed and cooperative management of virtual environments over a large-scale
distributed infrastructure.  Future work as well as opportunities are addressed in
Section~\ref{sec:future}. Finally Section~\ref{sec:conclusion} concludes this chapter.
